% =========================================================
% Proof: Power Rule (Special Case of Extended Product Rule)
% Source: volume-iii/analysis/differentiation/notes/notes-algebra.tex
% Mode: standard
% =========================================================

\subsection*{Power Rule (Special Case)}
\label{prf:power-rule-special-case}

\begin{remark*}[Return]
\hyperref[cor:power-rule-special-case]{$\leftarrow$ Back to Corollary (Power Rule, Special Case) in Notes}
\end{remark*}

\begin{proof}
Let $I \subseteq \mathbb{R}$ be an open interval, let $c \in I$, let
$f : I \to \mathbb{R}$ be differentiable at $c$, and let $n \in \mathbb{N}$.
The claim is that $f^n$ is differentiable at $c$ and
$(f^n)'(c) = n(f(c))^{n-1}f'(c)$.

\ProofStep{1}{Apply the Extended Product Rule with all factors equal to $f$.}
Set $f_1 = f_2 = \cdots = f_n = f$. By \ByThm{Extended Product Rule},
\[
  (f^n)'(c)
  = (f_1 f_2 \cdots f_n)'(c)
  = \sum_{k=1}^{n} f_k'(c) \prod_{\substack{i=1\\i\neq k}}^{n} f_i(c).
\]

\ProofStep{2}{Simplify each term.}
Since $f_k = f$ for every $k$, each factor $f_k'(c) = f'(c)$. The
remaining product contains $n-1$ factors each equal to $f(c)$:
\[
  \prod_{\substack{i=1\\i\neq k}}^{n} f_i(c)
  = (f(c))^{n-1}.
\]
Each term in the sum is therefore $f'(c) \cdot (f(c))^{n-1}$,
independently of $k$.

\ProofStep{3}{Evaluate the sum.}
Since the summand is constant in $k$, summing over $k = 1, \ldots, n$
gives a factor of $n$:
\[
  (f^n)'(c)
  = \sum_{k=1}^{n} f'(c)(f(c))^{n-1}
  = n \cdot f'(c) \cdot (f(c))^{n-1}
  = n(f(c))^{n-1}f'(c).
\]
\end{proof}

\begin{remark*}[Proof shape]
The argument is a direct specialisation of the Extended Product Rule:
set all $n$ factors equal to $f$, observe that every term in the sum
collapses to the same value $f'(c)(f(c))^{n-1}$, and count the $n$
identical terms. No induction is needed beyond what is already encoded
in the Extended Product Rule.

\FlashProof{1. Set $f_1=\cdots=f_n=f$ in the Extended Product Rule.
  2. Each term collapses to $f'(c)(f(c))^{n-1}$ since all factors are
  equal. 3. Sum of $n$ identical terms gives $n(f(c))^{n-1}f'(c)$.}
\end{remark*}

\begin{remark*}[Commentary]
\textbf{(a) Why does each product term equal $(f(c))^{n-1}$?}
In the Extended Product Rule, the $k$-th term differentiates factor
$f_k$ and evaluates all remaining $n-1$ factors at $c$. Since every
factor is $f$, the remaining product is $f(c)$ multiplied by itself
$n-1$ times, which is $(f(c))^{n-1}$. The index $k$ identifies which
factor is differentiated, but since they are all the same function,
the choice of $k$ does not change the result.

\textbf{(b) Why is summing $n$ identical terms valid here?}
The summand $f'(c)(f(c))^{n-1}$ depends on neither $k$ nor any
running variable — it is a fixed real number determined entirely by
$f$, $c$, and $n$. A sum of $n$ copies of a constant equals $n$
times that constant. This is not an approximation; it is exact
arithmetic.

\textbf{(c) What does this prove that the basic Product Rule does not?}
The basic Product Rule handles $n = 2$. The Extended Product Rule
extends this to arbitrary finite $n$ by induction. The Power Rule
Special Case then uses the Extended Product Rule to handle the
specific pattern where all $n$ factors are the same function, yielding
the clean formula $n(f(c))^{n-1}f'(c)$. Each result builds on the
previous one in a strict dependency chain.
\end{remark*}

\begin{remark*}[Dependencies]
\begin{itemize}
  \item \hyperref[cor:extended-product-rule]{Extended Product Rule}:
    The single result applied in Step 1. All remaining steps are
    arithmetic specialisation.
\end{itemize}

\FlashUses{cor:extended-product-rule}
\end{remark*}